% Generated draft: requires mathematical and source review.
\documentclass[11pt]{article}
\usepackage[margin=25mm]{geometry}
\usepackage{fontspec}
\setmainfont{Latin Modern Roman}
\usepackage{amsmath,amssymb,mathtools,unicode-math}
\setmathfont{Latin Modern Math}
\usepackage{enumitem,needspace,xurl,hyperref}
\setlength{\parindent}{0pt}
\setlength{\parskip}{5pt}
\setlength{\emergencystretch}{3em}
\newcommand{\N}{\mathbb{N}}
\newcommand{\Z}{\mathbb{Z}}
\newcommand{\Q}{\mathbb{Q}}
\newcommand{\R}{\mathbb{R}}
\newcommand{\C}{\mathbb{C}}
\newcommand{\F}{\mathbb{F}}
\newcommand{\A}{\mathbb{A}}
\newcommand{\PP}{\mathbb{P}}
\newcommand{\E}{\mathbb{E}}
\newcommand{\eps}{\varepsilon}
\newcommand{\dd}{\,\mathrm{d}}
\begin{document}
\section*{UnsolvedMath ID: 2307019}
\textbf{Source problem code:} catalogue entry
\subsection{Definitions and mathematical statement}
Let $Q=[0,1]^2$ and, for a set $\Omega$ of admissible integer pairs, put
\[
V_{\Omega}=\operatorname{span}\{(x,y)\mapsto x^my^n:(m,n)\in\Omega\},
\]
where span means finite linear combinations. In the conventional nonnegative-exponent setting these monomials are continuous. Density in $L^2(Q)$ means closure in the norm from two-dimensional Lebesgue measure; density in a continuous-function space means closure in the uniform norm.

\textbf{Problem (J. Korevaar).} Characterize the lattice-point sets $\Omega$ for which $V_{\Omega}$ is dense in $L^2(Q)$, or in the source's space $C_0$ on $Q$.

\textbf{Conjectured sufficient condition.} In the positive-exponent interpretation $\Omega\subset\mathbb Z_{>0}^2$, suppose that for some fixed $\varepsilon>0$ every $(m,n)\in\Omega$ satisfies
\[
\varepsilon m\leq n\leq m/\varepsilon,
\qquad\text{and}\qquad
\sum_{(m,n)\in\Omega}\frac1{mn}=\infty.
\]
The source conjectures that this suffices for the spanning property under consideration. The sum is over the set of lattice points, not an enumeration with arbitrary repetitions. The angle condition concerns exponent coordinates, not the variables of a monomial. The source attributes to Hellerstein that this divergence condition is not necessary.

The meaning of $C_0$ is essential and absent. In the usual vanish-at-infinity definition on a compact square it equals $C(Q)$, but positive-exponent monomials vanish on both coordinate edges and cannot approximate arbitrary boundary values. A possible intended space is
\[
\{u\in C(Q):u(x,0)=u(0,y)=0\text{ for all }x,y\in[0,1]\},
\]
but this is not silently substituted for the undefined notation. The general lattice domain is likewise unspecified: zero exponents require a convention for the displayed reciprocal sum, and negative exponents need not belong to the stated spaces. The supplied summary reports the question open as of the 2018 edition, with no definitive resolution located.
\subsection{Short English statement}
Characterize dense families of two-variable monomials. A reciprocal-exponent divergence condition is conjectured sufficient inside a fixed exponent cone, but the continuous target space and permitted lattice points need definition.
\subsection{Sources}
\begin{itemize}
\item Hayman, \emph{Research Problems in Function Theory} (2018), Problem 7.19, attributed to J. Korevaar, \url{https://arxiv.org/abs/1809.07200}; source-TeX extraction, accessed 2026-07-29. The 2018 update was retained for status triage.
\item Hellerstein, the reported nonnecessity of the reciprocal-sum condition; no publication, date, theorem location, or counterexample is supplied.
\item Local registry \texttt{data/unsolvedmath/problems.json}, permanent ID $2307019$, source code catalogue entry, \emph{Research Problems in Function Theory --- Problem 7.19}. The record labels this open and OPEN-TRIAGE, with NEEDS\_REVIEW source-status and rights qualifications.
\end{itemize}
\end{document}
